% ADDED figure (2026-11-02): the diagrams of one example box (C6a/C6b/§CA.7); generated by qft-C/tools/c6diag/figs/make_figs.sh
\tikzset{
  fdl/.style={line width=0.8pt, black},
  fdx/.style={circle, fill=black, inner sep=0pt, minimum size=3.6pt},
  fdv/.style={circle, fill=figB, draw=figB, inner sep=0pt, minimum size=5pt},
  fdlab/.style={font=\scriptsize, inner sep=2.6pt},
  fdmom/.style={->, >=Stealth, figO, line width=0.6pt},
}
\tikzset{
  fda/.style={fdl, postaction={decorate}, decoration={markings, mark=at position 0.55 with {\arrow{Stealth[length=5pt]}}}},
  fdS/.style={font=\small},
}
\begin{tikzpicture}[font=\small]
  \tikzset{fdo/.style={circle, draw=black, fill=white, line width=0.6pt, inner sep=0pt, minimum size=4.4pt}}
  \newcommand{\quadfig}[5]{%
    \begin{scope}[xshift=#1cm, yshift=#2cm]
      \coordinate (Q1) at (0,0.4); \coordinate (Q2) at (0.9,0.4); \coordinate (Q3) at (0,-0.4); \coordinate (Q4) at (0.9,-0.4);
      \foreach \a/\b in {#3} {\draw[fdl] (Q\a)--(Q\b);}
      \foreach \a/\b in {#3} {\node[fdx] at (Q\a){}; \node[fdx] at (Q\b){};}
      \foreach \c in {#4} {\node[fdo] at (Q\c){};}
      \node[fdlab,left] at ($(Q1)+(-0.06,0)$){$1$}; \node[fdlab,right] at ($(Q2)+(0.06,0)$){$2$};
      \node[fdlab,left] at ($(Q3)+(-0.06,0)$){$3$}; \node[fdlab,right] at ($(Q4)+(0.06,0)$){$4$};
      \node[fdS] at (0.45,-0.9) {#5};
    \end{scope}}
  \newcommand{\trifig}[4]{%
    \begin{scope}[xshift=#1cm, yshift=-5.95cm]
      \coordinate (Q1) at (0,0.4); \coordinate (Q2) at (0.9,0.4); \coordinate (Q3) at (0.45,-0.4);
      \foreach \a/\b in {#2} {\draw[fdl] (Q\a)--(Q\b);}
      \foreach \a/\b in {#2} {\node[fdx] at (Q\a){}; \node[fdx] at (Q\b){};}
      \foreach \c in {#3} {\node[fdo] at (Q\c){};}
      \node[fdlab,left] at ($(Q1)+(-0.06,0)$){$1$}; \node[fdlab,right] at ($(Q2)+(0.06,0)$){$2$};
      \node[fdlab,below] at ($(Q3)+(0,-0.04)$){$3$};
      \node[fdS] at (0.45,-1.15) {#4};
    \end{scope}}
  % (a) the seven terms with a normal-ordered remainder
  \node[fdS,anchor=west] at (-0.5,0.95) {(a) a field left inside the colons (open circle): $\langle0|\!:\!\cdots\!:\!|0\rangle=0$};
  \quadfig{0.00}{0}{}{1,2,3,4}{$0$}
  \quadfig{2.25}{0}{1/2}{3,4}{$0$}
  \quadfig{4.50}{0}{1/3}{2,4}{$0$}
  \quadfig{6.75}{0}{1/4}{2,3}{$0$}
  \quadfig{9.00}{0}{2/3}{1,4}{$0$}
  \quadfig{11.25}{0}{2/4}{1,3}{$0$}
  \quadfig{13.50}{0}{3/4}{1,2}{$0$}
  % (b) the three complete contractions
  \node[fdS,anchor=west] at (-0.5,-1.9) {(b) the complete contractions survive, $\langle0|0\rangle=1$:};
  \quadfig{1.6}{-2.85}{1/2,3/4}{}{$D_F(x_1-x_2)D_F(x_3-x_4)$}
  \quadfig{6.6}{-2.85}{1/3,2/4}{}{$D_F(x_1-x_3)D_F(x_2-x_4)$}
  \quadfig{11.6}{-2.85}{1/4,2/3}{}{$D_F(x_1-x_4)D_F(x_2-x_3)$}
  \node at (4.55,-2.85) {$+$}; \node at (9.55,-2.85) {$+$};
  % (c) three fields
  \node[fdS,anchor=west] at (-0.5,-5.0) {(c) three fields (Step 5): every term keeps a field, so $\langle0|T\{\hat\phi_1\hat\phi_2\hat\phi_3\}|0\rangle=0$};
  \trifig{2.4}{}{1,2,3}{$0$}
  \trifig{5.2}{1/2}{3}{$0$}
  \trifig{8.0}{1/3}{2}{$0$}
  \trifig{10.8}{2/3}{1}{$0$}
\end{tikzpicture}
